\cotefile{15}
\scanpage{172}
Calculs de cohomologie absolue des Motifs de Tate $\mathbb{Z}(n)$.

% archive sheet number: 171

\scanpage{173}
\margin{$\circ\!\!-\!\!-\quad \omega_\ell(q)={}$ordre de $q\bmod\ell$ dans $(\mathbb{Z}/\ell\mathbb{Z})^{*}$}

\diagram{\begin{tikzcd}[ampersand replacement=\&,column sep=large]
k=\mathbb{F}_q \arrow[r,dash,"\pi\simeq\widehat{\mathbb{Z}}"'] \& \bar{k}
\end{tikzcd}}
engendré par $f=\mathrm{frob}^{q}$\qquad $M$ un $\mathbb{Z}_\ell$-module de t.f. sur le quel $\pi$ opère continûment.
\[
\begin{aligned}
H^0(\pi,M)&=M^\pi=M^{(1-f)}=\mathrm{Ker}\,(1-f)_M\\
H^1(\pi,M)&=M_{\deleted{\unsure{f|i}}\pi}=M_{(1-f)}=\mathrm{Coker}\,(1-f)_M=M/(1-f)M\\
H^i(\pi,M)&=0\qquad\text{si }\ i\geq 2
\end{aligned}
\]

\[
H^0(\pi,\mathbb{Z}_\ell(n))=0\quad\text{si }\ n\ne0,\quad =\mathbb{Z}_\ell\quad\text{si }\ n=0
\]
\[
\begin{aligned}
H^1(\pi,\mathbb{Z}_\ell(n))&\underset{\text{\inserted{non canoniquement}}}{=}\mathbb{Z}_\ell/(1-q^n)\mathbb{Z}_\ell=\begin{cases}
\mathbb{Z}_\ell&\text{si }\ n=0\\
0&\text{si }\ \ell\nmid(1-q^n)\quad\text{i.e. }\deleted{\unsure{n\ne0}}\ \inserted{n\not\equiv0\ (\omega_\ell(q))}\\
\mathbb{Z}/\ell^{\unsure{v|\nu}_\ell(1-q^n)}&\text{si }\ \underbrace{\ell\mid1-q^n}_{\inserted{\text{i.e. }n\equiv0\ (\omega_\ell(q))}},\ n\ne0
\end{cases}
\end{aligned}
\]
\margin{$\circ\!\!-\!\!-$}
\[
\left\{\begin{aligned}
H^{i}(\pi,\mathbb{Q}_\ell(n))&=0\quad\text{si }\ n\ne0,\quad\text{quel que soit }i\\
H^0(\pi,\mathbb{Q}_\ell)&=\mathbb{Q}_\ell\deleted{\unsure{f|1}},\ H^1(\pi,\mathbb{Q}_\ell)=\mathbb{Q}_\ell
\end{aligned}\right.
\]

\[
\deleted{H^i(\pi,\underline{\mathbb{Z}}(n))}
\]
\[
\left\{\begin{aligned}
\mathbb{H}^i(\mathbb{F}_q,\underset{\sim}{\mathbb{Q}}(n))&=0\quad\text{si }\ n\ne0,\quad\text{quel que soit }i\\
\mathbb{H}^0(\mathbb{F}_q,\underset{\sim}{\mathbb{Q}})&=\mathbb{Q},\ \mathbb{H}^1(\mathbb{F}_q,\underset{\sim}{\mathbb{Q}})=\underset{\sim}{\mathbb{Q}}.
\end{aligned}\right.
\]
\[
\left\{\begin{aligned}
H^0(\mathbb{F}_q,\underset{\sim}{\deleted{\mathbb{Q}}\mathbb{Z}})&=\mathbb{Z}\unsure{,|{}_\ell}\ (\text{mod }p\text{-gpes }?),\quad H^1(\mathbb{F}_q,\underset{\sim}{\mathbb{Z}})=\mathbb{Z}\quad(\text{mod }p\text{-gpes }?)\\
H^0(\mathbb{F}_q,\underset{\sim}{\mathbb{Z}}(\underset{\substack{\vert\\\inserted{n\ne0}}}{n}))&=0\quad(\text{id}),\quad H^1(\mathbb{F}_q,\underset{\sim}{\mathbb{Z}}(\underset{\substack{\vert\\\inserted{n\ne0}}}{n}))=\mathbb{Z}/(1-q^n)\mathbb{Z}\quad(\text{mod }p\text{-gpes }?)
\end{aligned}\right.
\]

\noindent\rule{.76\linewidth}{.35pt}
\[
\begin{aligned}
\mathbb{H}^i(\mathrm{Spec}\,\mathbb{Z},\mathbb{Z}(1))&=\begin{cases}0&\text{\unsure{si|sf}}\deleted{\unsure{<}}\ i\ne3\\\mathbb{Z}&\text{si }\ i=3\end{cases}\\
\mathbb{H}^i(\mathrm{Spec}\,\mathbb{Z},\mathbb{Z})&=\begin{cases}0&\text{si }\ i\ne0\\\mathbb{Z}&\text{si }\ i=0\end{cases}
\end{aligned}
\]
\[
H^i(\mathrm{Spec}(\mathbb{Z}),\deleted{\unsure{\mathbb{Z}}}\mathbb{Q}(n))\xrightarrow{\sim} H^{\unsure{i|1}}(\deleted{\unsure{k}}U,\deleted{\unsure{\mathbb{Z}}}\mathbb{Q}(n))
\]
si $U$ ouvert de Zariski de $\mathrm{Spec}\,\mathbb{Z}$, et $n\ne0$ (tout $i$).

% archive sheet number: 172

\scanpage{174}
$\mathrm{Spec}\,\mathbb{Z}$

\figure{\begin{tikzpicture}
\draw (0,0) .. controls (.8,-.02) and (1.4,.04) .. (2,0.04) .. controls (3,.05) and (3.55,-.08) .. (4.4,-.02);
\draw[gray!45,thin] (1.46,.37)--(1.55,.39);
\fill (1.8,.035) circle (1.5pt) node[below] {$x$};
\draw[thin] (4.08,-.05)--(4.05,-.12);
\end{tikzpicture}}
\[
i_x^!(M)=\unsure{1|\vert}\,i_x^*(M)(-1)[-2]
\]
\[
\cdots\unsure{-|\to}H^{i-2}(x,i_x^*(M)(-1))\longrightarrow H^i(X,M)\unsure{-|\to}H^{\unsure{i|2}}(U,M)\unsure{-|\to}H^{i-1}(x,i_x^*(M)(-1))\to\cdots
\]
\[
M=\underset{\sim}{\deleted{\mathbb{Z}}\mathbb{Q}}(n)
\]
\margin{\unsure{peu|pour} de mérite, car si $i\geq4$, les deux mbs doivent être $=0$}
\[
\deleted{H^i}\quad
\left\{\begin{aligned}
i\ne1,2,3\qquad&H^i(X,M)\xrightarrow{\sim}H^i(U,M)\\
0\to H^1(X,M)&\to H^1(U,M)\to H^0(S,i_x^*(M)(-1))\longrightarrow H^2(X,M)\\
&\longrightarrow H^2(U,M)\to H^1(S,i_{\unsure{x|\bullet}}^*(M)(-1))\to H^3(X,M)\\
&\longrightarrow\underset{\inserted{=0\,?}}{H^3(U,M)}\to\deleted{\unsure{H}}0
\end{aligned}\right.
\]

\[
M=\deleted{\mathbb{Z}}\mathbb{Q}_{\deleted{\unsure{S|\ell}}}(n),\quad i_{\unsure{\bullet|x}}^*(M)(\unsure{-1})\deleted{\unsure{x}}=\mathbb{Q}(n-1)_S
\]
\[
\text{\uline{Si $n\ne1$}}\qquad H^i(X,\mathbb{Q}(n))\xrightarrow{\sim}H^i(U,\mathbb{Q}(n))\qquad\text{\uline{tout $i$}}
\]
\margin{[est-ce nul $??$]}

Si $n=1$
\[
\begin{aligned}
0\unsure{-|\to}H^1(X,\mathbb{Q}(1))&\to H^1(U,\mathbb{Q}(1))\to\deleted{\unsure{\mathbb{Z}}}\mathbb{Q}^S\longrightarrow H^2(X,\mathbb{Q}(1))\to\\
&H^2(U,\mathbb{Q}(1))\to\mathbb{Q}^S\longrightarrow H^3(X,\mathbb{Q}(1))\longrightarrow0
\end{aligned}
\]
$H^i(X,\mathbb{Q}(1))=0$ sauf pour $i=3$, où on trouve $\mathbb{Q}$ !

% archive sheet number: 173

\scanpage{175}
% Three long oblique strokes run through the working on this sheet.
\deleted{
\diagram{\begin{tikzcd}[ampersand replacement=\&]
k=\mathbb{F}_q \arrow[r,dash,"\pi=\widehat{\mathbb{Z}}"'] \& \bar{k}
\end{tikzcd}}
engendré par $\unsure{\mathrm{Frob}|\mathrm{frob}}^{q}_{\deleted{\unsure{q}}}$

\[
\begin{aligned}
H^0(\pi,M)&=M^\pi\\
H^1(\pi,M)&=\varprojlim_n H^1\underbrace{(\pi,M/\ell^n M)}=
\end{aligned}
\]

a) si $\pi$ opère trivialement, c'est $\mathrm{Hom}(\pi,M)$

b) Si $M$ est fini, \unsure{donc} $\pi$ opère à travers $\pi/U\simeq\mathbb{Z}/n\mathbb{Z}$

\margin{$H^1(\pi,M)$}
\diagram{\begin{tikzcd}[ampersand replacement=\&,column sep=small,row sep=small]
0\unsure{-|\to} H^1(\pi/U,\underset{\substack{\shortparallel\\\inserted{M}}}{M^U})\unsure{-|\to} H^1(\pi,M)\arrow[r] \&
\underset{\substack{\shortparallel\\\inserted{\mathrm{Hom}(U,M^\pi)}}}{H^0(\pi/U,\mathrm{Hom}(U,M))}\arrow[d]\\
\& H^2(\pi/U,\underset{\substack{\shortparallel\\\inserted{M}}}{M^U})
\end{tikzcd}}
\[
\deleted{M_\pi}\quad\frac{\mathrm{Ker}\,N}{\mathrm{Im}\,(1-f)}
\]

\[
\begin{aligned}
\varphi(xy)&=\varphi(x)+x\varphi(y)\\
\varphi(x)&=xa-a\\
\varphi(xy)&=xya-a\\
&=xa-a+x(ya-a)\\
&=\deleted{\unsure{xya-a}}
\end{aligned}
\]
\[
\underset{\substack{\shortparallel\\1+f+\cdots+f^{n-1}}}{N}\qquad 1-f
\]
\[
\varphi(f^n)=
\begin{cases}
(1+x+\cdots+x^{n-1})\varphi(x)=\dfrac{1\deleted{\unsure{+}}-x^n}{1-x}\,\varphi(x)&\text{si }n\geq1\\
0&\text{si }n=0\\
\unsure{(}&
\end{cases}
\]
\[
\varphi(f^{\unsure{n|m}})=1+x+\phantom{x} +x^{n-1}
\]
\[
\begin{aligned}
&\deleted{\varphi(x^n)=\varphi(x)+}\\
\varphi(x^{\unsure{2|1}})&=\varphi(x)+x\varphi(x)\\
\varphi(x^3)&=\varphi(x)+x\varphi(x)\unsure{-|+}x^{\unsure{2|3}}\varphi(x)\\
\varphi(x^n)&=(1+x+\cdots+x^{n-1}\deleted{\unsure{(x)}})\varphi(x)\\
0=\varphi(x^{-1}x)&=\varphi(x^{-1})+x^{-1}\varphi(x)\\
\deleted{\unsure{\varphi(x}}\qquad\varphi(x^{-1})&=-x^{-1}\varphi(x)\\
\varphi(x^{-n})&=-(1+x^{-1}+\cdots+x^{-(n-1)})x^{-1}\varphi(x)\\
&\deleted{\unsure{=-(x}}\quad-\frac{1-x^{-n}}{1-x^{-1}}x^{-1}\varphi(x)\\
&=\frac{1-x^{-n}}{1-x}
\end{aligned}
\]

\[
\begin{aligned}
\unsure{\varphi(f)}\qquad Z^1(\mathbb{Z},M)&\simeq M\\
B^1(\mathbb{Z},M)&\simeq (1-f)M\\
H^1(\mathbb{Z},M)&\simeq \deleted{\unsure{1}}M/(\deleted{\unsure{1}}1-f)M
\end{aligned}
\]
\[
\frac{1-f^n}{1-f}\,\varphi(f)
\]
Si $n\equiv0\quad(N)$
\[
1+f+\cdots+f^{n-1}
\]
}

% archive sheet number: 174

\scanpage{176}
Lorsque \deleted{$X$ est \unsure{l'on|on} est sous les conditions de} les anneaux localisés stricts de $X$ sont factoriels, on tire aisément de (5.3.)
\[
H^1(X,\underline{\unsure{\mathrm{Div}|\mathrm{Dv}|\mathcal{D}\mathrm{iv}}}_X)=0
\]
d'où \deleted{la proposition}

\uline{Proposition} 5.\unsure{4}. Supposons que l'on soit sous les conditions de 5.1. a), \inserted{$X$ étant irréductible de pt gén. $y$.} Alors les homomorphismes canoniques
\[
H^2(X,\mathrm{G}_{m\,X})\unsure{-|\to}H^2(X,\mathrm{R}_X^*)\longrightarrow H^2(y,\mathrm{G}_{m\,y})=\mathrm{Br}(y)
\]
sont injectifs.

Comme $\mathrm{Br}(X)\subset H^2(X,\mathrm{G}_{m\,X})$, il s'ensuit

\uline{Corollaire} 5.5. Sous les \unsure{mêmes} conditions,
\[
\mathrm{Br}(X)\unsure{-|\to}\mathrm{Br}(y)
\]
est injectif.

Ainsi, $\mathrm{Br}(X)$ s'identifie \inserted{alors} à un \unsure{sous-gpe|sous-groupe} du groupe de Brauer $\deleted{\unsure{\mathrm{B}}}\mathrm{Br}(K)$ du corps des fonctions $K=k(y)$ de $X$.

\uline{Théorème} 5.6. \inserted{Soit $X$ un préschéma, $\unsure{\zeta|\xi}\in H^2(X,\mathrm{G}_m)$} a) Si $X$ est un schéma local, et \deleted{$\zeta\in H^2(X,\mathrm{G}_{m\,X})$}, pour que $\zeta$ \deleted{\unsure{provient|provienne} d'un élément de} \inserted{soit dans $\mathrm{Br}(X)$, il} faut et s. qu'il existe un \deleted{extension} morphisme fini étale surjectif (galoisien si on veut) $f\colon X'\to X$, \deleted{tel que $\zeta|X'$ soit} qui \deleted{\unsure{splitt|splitte} $\zeta$} annule $\zeta$.

% archive sheet number: 175

\scanpage{177}
\[
\deleted{\mu_n\longrightarrow\mathrm{G}_m\xrightarrow{n}\mathrm{G}_m}
\]

$\mu$

\[
\deleted{H^3(X,\mathrm{G}_m)\longrightarrow H^3}
\]
\[
\deleted{H_x^i(}
\]
\[
H^i_S(X,\mathrm{G}_m)\unsure{-|\to}H^{\deleted{\unsure{0|3}}i}(X,\mathrm{G}_m)\to H^{\deleted{\unsure{0|3}}i}(U,\mathrm{G}_m)\unsure{-|\to}H^{\unsure{i+1|i-1}}_S(X,\mathrm{G}_m)\to
\]
\[
\underline{H}^i_x(\mathrm{G}_m)=
\begin{cases}
0&\text{si }i=0\\
\mathbb{Z}_x&\text{si }i=1\\
0&\text{si }i\geq2
\end{cases}
\]
\[
H^i_S(X,\mathrm{G}_m)=H^{i-1}(S,\mathbb{Z})=
\left\{\begin{array}{lll}
\deleted{\mathbb{Z}}0&\text{si }i=0\\
\mathbb{Z}^S&\text{si }i=1\\
0&\text{si }i=2&\smash{\left.\rule[-2.3ex]{0pt}{6.5ex}\right]}\\
H^1(S,\mathbb{Q}/\mathbb{Z})&\text{si }i=3\\
0&\text{si }i\geq4
\end{array}\right.

\]
\[
0\unsure{-|\to}\mathbb{Z}\to\mathbb{Q}\unsure{-|\to}\mathbb{Q}/\mathbb{Z}\unsure{-|\to}0
\]

\[
\begin{aligned} i&=1,3\\i&=0,2\end{aligned}
\qquad\big/\qquad \text{Si }i\ne0,1,2,3,\qquad H^i(\bar{X},\mathrm{G}_m)\xrightarrow{\sim}H^i(U,\mathrm{G}_m)
\]
\margin{peu de mérite, car les deux sont nuls.}

\[
\begin{aligned}
0&\to\underset{\substack{\shortparallel\\\inserted{\mathbb{Z}^*=\{+1,-1\}}}}{H^0(\bar{X},\mathrm{G}_m)}
\to H^0(\bar{U},\mathrm{G}_m)\to\mathbb{Z}^S\to
\underset{\substack{\shortparallel\\\inserted{0}}}{H^1(\bar{X},\mathrm{G}_m)}
\unsure{-|\to}\underset{\substack{\shortparallel\\\inserted{0}}}{H^1(\bar{U},\mathrm{G}_m)}
\to\underset{\substack{\shortparallel\\\inserted{0}}}{H^2_S(\bar{X},\mathrm{G}_m)}
\end{aligned}
\]
\[
\begin{aligned}
0&\to\underset{\substack{\shortparallel\\\inserted{0}}}{H^{\unsure{2|1}}(\bar{\deleted{\unsure{U|X}}X},\mathrm{G}_m)}
\to H^2(\bar{U},\mathrm{G}_m)\to\underset{\substack{\shortparallel\\\inserted{(\mathbb{Q}/\mathbb{Z})^S}}}{H^1(S,\mathbb{Q}/\mathbb{Z})}
\xrightarrow[\inserted{\substack{\text{\uline{somme}}\\\text{\uline{des inv.}}}}]{}
\underset{\substack{\shortparallel\\\inserted{\mathbb{Q}/\mathbb{Z}}}}{H^3(\bar{X},\mathrm{G}_m)}
\to\underset{\substack{\shortparallel\\\inserted{0\,?}}}{H^3(\bar{U},\mathrm{G}_m)}\to0
\end{aligned}
\]

% archive sheet number: 176

\scanpage{178}
\[
\begin{aligned}
H^0(\bar U,\mathrm{G}_m)&=\mathbb{Z}/2\mathbb{Z}+\mathbb{Z}^S\\
H^1(\bar U,\mathrm{G}_m)&=0\\
H^2(\bar U,\mathrm{G}_m)&=\mathrm{Ker}\left((\mathbb{Q}/\mathbb{Z})^S\longrightarrow\mathbb{Q}/\mathbb{Z}\right)\\
H^i(\bar U,\mathrm{G}_m)&=0\qquad\text{si }i\geq3
\end{aligned}
\]
\[
0\to\mu_{\ell^n}\to\mathrm{G}_m\xrightarrow{\ell^n}\mathrm{G}_m\to0
\]
\deleted{\unsure{alors|d'où}}
\[
\begin{aligned}
H^1(\bar U,\mu_{\ell^n})&\simeq(\mathbb{Z}/\ell^n\mathbb{Z})^S\\
H^2(\bar U,\mu_{\ell^n})&\simeq\mathrm{Ker}\left((\mathbb{Z}/\ell^n\mathbb{Z})^S\longrightarrow(\mathbb{Z}/\ell^n\mathbb{Z})\right)\\
H^i(\bar U,\mu_{\ell^n})&=0\qquad\text{si }i\geq3
\end{aligned}
\]
\[
H^0(\bar U,\mathbb{Z}_\ell)\simeq\mathbb{Z}_\ell
\]
\[
\left\{\begin{aligned}
H^1(\bar U,\mathbb{Z}_\ell(1))&\simeq\mathbb{Z}_\ell^S\\
H^2(\bar U,\mathbb{Z}_\ell(1))&\simeq\mathrm{Ker}\left(\mathbb{Z}_\ell^S\longrightarrow\mathbb{Z}_\ell\right)
\end{aligned}\right.
\]
\margin{petite modification pour $\ell=2$ en dim $0$}
\[
\left\{\begin{aligned}
H^0(X,\mathbb{Q}(1))&=\deleted{\unsure{\mathbb{Q}}}0\\
H^1(X,\mathbb{Q}_{\unsure{\ell|\vphantom{\ell}}}(1))&=\deleted{\unsure{0}}0\\
H^2(X,\mathbb{Q}_{\unsure{\ell|\vphantom{\ell}}}(1))&=0\\
H^3(X,\mathbb{Q}(1))&=\deleted{\unsure{0}}\mathbb{Q}
\end{aligned}\right.
\]

% archive sheet number: 177

\scanpage{179}
% The opening passage is enclosed and crossed by long diagonals.
\deleted{On peut supposer $X$ irréductible de pt générique $y$. L'inclusion $y\unsure{-|\to}X$ définit une suite spectrale (pour la cohomologie étale)
\[
H^*(y,\mathrm{G}_m)\Longleftarrow E_2^{pq}=H^p(X,\mathbb{R}^qi_{\unsure{*|x}}(\mathrm{G}_m))
\]
Or \deleted{\unsure{comme les}} comme les $H^q(y,\mathrm{G}_m)$ et les fibres de $\mathrm{R}^q i_*(\mathrm{G}_m)$, qui s'explicitent par de la cohomologie galoisienne, sont \deleted{\unsure{des}} groupes de torsion pour $q>0$, on trouve que \deleted{\unsure{l'on a} \unsure{sur|dans} les \unsure{les} groupes}
\[
H^p(X,\mathrm{R}_X^*)\tag*{$(*)$}
\]
\deleted{, où $\mathrm{R}_X=\mathrm{R}i_*(\mathrm{G}_m)$}\inserted{est un groupe de torsion pour $p>0$,} où $\mathrm{R}_X^*=\deleted{\unsure{\mathrm{R}i_*}}i_*(\mathrm{G}_{m\,\unsure{x|y}})$ est le faisceau (étale) des fonctions rationnelles \inserted{non nulles} sur $X$. \deleted{(Ceci étant}}

D'autre part, \inserted{si $X$ est \inserted{\unsure{noethérien}} ``sans idéaux premiers immergés'',} on a la suite exacte
\[
0\unsure{-|\to}\mathrm{G}_{m\,X}\longrightarrow\mathrm{R}_X^*\longrightarrow\underline{\mathrm{Div}}_X\longrightarrow0\tag*{(5.\unsure{1|2}.)}
\]
où $\underline{\mathrm{Div}}_X$ est le faisceau \inserted{(pour la \unsure{top|top.}} étale) des diviseurs (de Cartier) sur $X$, ce qui \deleted{\unsure{montre que}} \inserted{implique} compte tenu de $(*)$ que
\[
H^{i}(X,\underline{\mathrm{Div}}_X)\xrightarrow{\partial}H^{i+1}(X,\deleted{\unsure{\mathrm{R}|\mathrm{D}}}\mathrm{G}_{m\,X}).
\]
% slanted
\margin{\uline{donner le} \uline{corollaire} 5.4.}
est bijectif mod torsion \deleted{\unsure{si}} si $i\geq1$. \deleted{\unsure{Ainsi} \unsure{on aura|ne sera pas} vraiment utilisé} L'hypothèse de factorialité d'autre part implique \deleted{l'\unsure{isomorphisme}} qu'on a un isomorphisme de faisceaux étales :

% archive sheet number: 178

\scanpage{180}
Calculs de cohomol. absolue du motif de Tate
\[
\unsure{\mathbb{Z}|N}(n).
\]

% archive sheet number: 179

\scanpage{182}
Généralités fonctoriello-catégoriques sur la distribution des frobenius dans les gpes de Galois motiviques.

% archive sheet number: 181

\scanpage{183}
Soit $\mathcal M$ une catégorie de motifs géométriques, \inserted{de type \unsure{lié|lisse|fini},} sur le préschéma $S$ de type fini sur $\mathbb Z$. Pour simplifier, supposons que les degrés de $\mathcal M$ soient pairs (sinon, il sera possible, pour \deleted{la plupart des} \inserted{certaines} questions de distribution des Frobenius, de remplacer $\mathcal M$ par $\mathcal M^{\mathrm{pair}}$). Tout choix d'un pt fermé $s\in S$ permet donc, de définir des (une \unsure{collection|notion} de) ``$\mathbb R$-points Motiviques'' \deleted{\unsure{point de vue|de $\mathcal M$}''} \inserted{de $\mathcal M$}.

Pour fixer les idées, choisissons un pt $\mathbb R$-motivique de $\mathcal M$, d'où un gpe \unsure{$\Lambda$physique|physique|algébrique} $G_{\mathbb R}$ sur $\mathbb R$.

Il y a lieu de \deleted{\unsure{dupliquer|d'expliciter}} \inserted{démultiplier} les \unsure{point|points} \inserted{fermés $s$} de $S$, en formant les couples $s'=(s,\unsure{i'|i|i_s|\xi})$, où $s\in S^0$, et
% slanted
\margin{l'exposant $0$ désigne \unsure{des|les} pts \uline{fermés}}
$i_{\deleted{\unsure{s|x}}}\colon G_{m\,\unsure{\mathbb C|C}}-G(s)_{\mathbb C}$ est de type \deleted{\unsure{désigné|divisé}} préférentiel --- $G(s)$ désigne le gpe de Galois motivique associé à $\mathcal M_s$.

\deleted{Alors} Soit $S'$ l'ens de ces couples. A tout $s'\in S'$ est associé un $\mathbb R$-point motivique de $\mathcal M$, \deleted{d'où un} déterminé à isom près, d'où un $\unsure{\zeta|\xi}(s')\in H^1(\mathbb R,\deleted{\unsure{F|G}}G_{\mathbb R}^{\mathrm{or}})$, où $G^{\mathrm{or}}$ est le noyau de $G_{\mathbb R}^0\xrightarrow{\varepsilon}G_m$. De plus, si $G_{\mathbb R}^{\zeta(s')}$ est ``le'' gpe déduit de $G_{\mathbb R}$ par torsion par $\zeta(s')$, on obtient un hom
\[
u(s')\colon G(s')_{\deleted{\unsure{R|k}}\mathbb R}=G(s)_{\mathbb R}\longrightarrow G_{\mathbb R}^{\zeta(s')}
\]
compatible avec $\varepsilon_{\mathbb R},J_{\mathbb R}$, et déterminé en fait mod \deleted{\unsure{automorphisme|automorphismes}} \inserted{conjugaison} \deleted{\unsure{intérieur|interne|intérieurs}} par un élément de $G_{\mathbb R}^{\mathrm{or}}(\mathbb R)$.
% archive sheet number: 182
